\chapter{Data Freshness Model}

\section{Design Principle}
Every API response includes freshness metadata so consumers can assess
data currency.

\section{Cache-Control Strategy}

\begin{table}[h]
\centering
\begin{tabular}{@{}lll@{}}
\toprule
\textbf{Resource} & \textbf{TTL} & \textbf{Invalidation trigger} \\
\midrule
Symbol metadata     & 1\,h   & File hash change \\
Dependency subgraph & 1\,h   & Source/target symbol change \\
Search results      & 15\,min & Any symbol mutation \\
Module import graph & 1\,h   & File dependency change \\
Materialised views  & 1\,h   & Cron refresh \\
\bottomrule
\end{tabular}
\end{table}

\section{HTTP Cache Headers}
Query API responses include:
\begin{lstlisting}[language={}]
Cache-Control: public, max-age=60, stale-while-revalidate=300
ETag: "sha256:<hash>"
\end{lstlisting}

\section{SSE Invalidation}
An optional \texttt{/api/v1/events/invalidation} SSE endpoint streams
real-time invalidation events. The frontend listens and calls
\texttt{queryClient.invalidateQueries} on receipt.

\section{Staleness Alerting}
\begin{itemize}
  \item \texttt{indexed\_at} $>$ 24\,h: warning badge in UI.
  \item \texttt{indexed\_at} $>$ 48\,h: Prometheus alert \texttt{index\_stale}.
  \item Redis unavailable: fallback to DB; log \texttt{cache\_miss\_fallback}.
\end{itemize}

\section{Consistency Guarantees}
Serialisable isolation for index writes (strong within transaction).
Cache reads are eventually consistent (worst-case staleness = TTL).
SSE reduces effective staleness to sub-second for connected clients.
